%!TeX program = xelatex
\documentclass[12pt,hyperref,a4paper,UTF8]{ctexart}
\usepackage{BNUpapers}

\title{\Huge\textbf{示例：城市绿地与夏季气温}}
\author{%
\kaishu\Large
    姓名\ \underline{张三} \qquad
    学号\ \underline{2024XXXXXX} \qquad
    院系\ \underline{某某学院}%
}
\date{\today}

%%-------------------------------正文开始---------------------------%%
\begin{document}

%%-----------------------封面--------------------%%
\cover

%%------------------摘要-------------%%
\begin{abstract}
本文以某市若干公园与街道为示例，说明课程论文中摘要、正文、图表和参考文献的写法。示例比较了绿地覆盖较高的街区与绿地较少的街区在夏季午后的气温差异，并给出一个简单的数量关系。结论仅用于演示版式，不对应真实观测。
\end{abstract}

\thispagestyle{empty} % 首页不显示页码

%%--------------------------目录页------------------------%%
\newpage
\tableofcontents

%%------------------------正文页从这里开始-------------------%
\newpage

%%可选择这里也放一个标题
%\begin{center}
%    \title{ \Huge \textbf{{标题}}}
%\end{center}

\section{研究说明}
本文是课程论文的示例正文。讨论对象是城市中的公共绿地与夏季午后气温，材料为虚构的街区记录，用来占位，方便替换成自己的题目。

示例按“问题、材料、结果、讨论”组织。编译方式见仓库说明：\verb|xelatex -> bibtex -> xelatex*2|。

本示例用到的材料包括：
\begin{itemize}
    \item 四条样例街道的绿地比例
    \item 对应的午后气温记录
    \item 一幅校名图，仅作插图演示
    \item 两则与科学史有关的参考文献
\end{itemize}

把标题、姓名、学号和院系改成自己的信息后即可继续写。封面、目录和参考文献页保留在示例中，页码从目录页开始。

\section{结果与讨论}
\subsection{一个简单关系}
样例里把温差写成行内式 $v-\varepsilon+\phi=2$。

若把同一关系单独列出，如\autoref{Euler}：
\begin{equation}
    v-\varepsilon+\phi=2
    \label{Euler}
\end{equation}

\subsection{插图}
校名图如\autoref{bnu}所示。\verb|~\autoref{}| 会生成“图”“式”等前缀。

\begin{figure}[!htbp]
    \centering
    \includegraphics[width =.4\textwidth]{figures/BNU_name.pdf}
    \caption{示例插图}
    \label{bnu}
\end{figure}

插图示例代码：

\begin{verbatim}
    \begin{figure}[!htbp]
        \centering
        \includegraphics[width =.4\textwidth]{figures/BNU_name.pdf}
        \caption{示例插图}
        \label{bnu}
    \end{figure}
\end{verbatim}

\subsection{补充说明}
需要强调的句子可以放进圆角文本框，命令是\verb|\tbox{}|。

\tbox{
    示例：绿地比例更高的街区，午后气温略低。
}

\subsection{样例数据}
四条街道的记录如\autoref{doc}所示。
\begin{table}[!htbp]
    \centering
    \begin{tabular}{l  | l}
    \hline
        街道 & 绿地比例 \\
        \hline
        甲 & 高 \\
        乙 & 中 \\
        丙 & 低 \\
        丁 & 很低 \\
        \hline
    \end{tabular}
    \caption{示例表格}
    \label{doc}
\end{table}

%\section{定理环境}
%\begin{Theorem}
%\end{Theorem}
%
%\begin{Lemma}
%\end{Lemma}
%
%\begin{Corollary}
%\end{Corollary}
%
%\begin{Proposition}
%\end{Proposition}
%
%\begin{Definition}
%\end{Definition}
%
%\begin{Example}
%\end{Example}
%
%\begin{proof}
%\end{proof}

\subsection{参考文献}
正文中用\verb|\cite{}|引用。

例如：


   \textit{ 这里引用了\cite{0Isaac}，也引用了\cite{2016The}。}


引用过的条目会出现在文末。

\section{小结}
\subsection{可以替换的部分}
\begin{itemize}
    \item 标题、摘要和各节文字
    \item 图、表和公式中的示例数据
    \item \texttt{reference.bib} 中的文献条目
\end{itemize}
\subsection{写作时注意}
示例只说明版式。正式作业应换成自己的问题、数据和文献，并核对学校对封面信息的要求。

\tbox{
把本页的示例句子换成自己的结论即可。
}
%%----------- 参考文献 -------------------%%
%在reference.bib文件中填写参考文献，此处自动生成

\reference


\end{document}